\section*{Open Questions --- Five Concrete Follow-ups}

\begin{enumerate}

\item \textbf{Extension to doped and multi-band Hubbard.}
\textit{Question:} How do the per-block HVA gate counts scale when the 2D Hubbard model is
extended to the doped regime (away from half-filling) or to the three-band Emery model
relevant for cuprate superconductivity? \\
\textit{Basis:} Cai's closed-form formulas $N_{1q}=4V^{3/2}+7V-4\sqrt{V}$ and
$N_{2q}=8V^{3/2}+V-4\sqrt{V}$ are derived for the single-band Hubbard model at generic
filling. The physics of interest lives at finite doping and often in multi-band models;
whether the $O(V^{3/2})$ scaling and constant prefactors survive is not addressed. \\
\textit{Next steps:} Implement the doped/three-band Hamiltonian in OpenFermion,
Jordan--Wigner transform, count HVA-block terms, and fit the empirical scaling.
Report new closed-form or empirical constants. Free endpoint: local OpenFermion + Cirq
at $V\le 12$.

\item \textbf{Tensor-hypercontraction / low-rank interaction factorisation.}
\textit{Question:} How much does modern tensor-hypercontraction (THC) or low-rank
factorisation of the two-body integrals reduce the two-qubit gate count vs.\ the sparse
fermionic-swap-network construction Cai uses? \\
\textit{Basis:} Kivlichan et al.'s swap network is optimal for arbitrary-connectivity
sparse Hamiltonians but suboptimal for structured (nearest-neighbour Hubbard) or low-rank
Hamiltonians. For chemistry, THC has driven $\sim 10\times$ reductions in two-qubit cost. \\
\textit{Next steps:} For $V\in\{9,16,25\}$, factor the interaction as a sum of rank-1
terms via OpenFermion's low-rank module, count Trotter-block gates, compare to Cai's
Appendix A2.

\item \textbf{Quantum--classical embedding (DMFT / CDMFT).}
\textit{Question:} How does the HVA resource estimate change when embedded in a DMFT
scheme where only the impurity cluster is quantum? \\
\textit{Basis:} Cai treats the full $5\times 5$ lattice as quantum; DMFT/CDMFT solves a
small impurity + self-consistent bath, changing both $V$ and boundary conditions and
therefore the qubit count and per-block gate count. \\
\textit{Next steps:} Set up a 4-site CDMFT impurity with a classical bath in OpenFermion,
Jordan--Wigner it, evaluate Cai's formulas at $V_{\mathrm{impurity}}$, and compare the
resulting resource estimate to the 50-qubit full-lattice number.

\item \textbf{Noise-adaptive resource estimation.}
\textit{Question:} How do the gate-count and error-budget estimates change under
noise-adaptive compilation (Pauli-frame randomisation, dynamical decoupling,
hardware-aware scheduling) rather than the uniform-error assumption? \\
\textit{Basis:} Cai's error budget assumes each two-qubit gate independently contributes
$p_{2q}=10^{-4}$ with no cross-talk, no idle-qubit decoherence, and no correlated error.
Real devices deviate; adaptive compilation partially mitigates. \\
\textit{Next steps:} Instantiate $V=4$ HVA in Cirq/Qiskit with a realistic noise model
(T1/T2 + depolarising + measurement), Monte-Carlo the mean per-shot error with and
without DD insertions, compare to the 2.5-error prediction.

\item \textbf{Classical-solver benchmark at matched target accuracy.}
\textit{Question:} At the same target accuracy (chemical accuracy per site, or
$\sim 10^{-3}$ of correlation energy), how does the projected quantum-resource cost
of the HVA-VQE compare to state-of-the-art AFQMC/DQMC (at accessible $U/t$) or DMRG
on the equivalent cylinder? \\
\textit{Basis:} Cai reports an isolated quantum estimate but does not benchmark it
against classical wall-clock cost at the same target accuracy on the same instance. \\
\textit{Next steps:} (i) Tabulate published DMRG/QMC $V{=}25$ energies + reported CPU-hour
cost at chemical accuracy; (ii) convert Cai's per-shot cost ($250\,\mu\mathrm{s} \times
N_{\mathrm{shots}}$) to wall-clock given a target variance; (iii) plot side-by-side.
Free endpoint: literature review only.

\end{enumerate}
